\chapter{Knot invariants and the Kontsevich integral}
\label{ch:kontsevich-integral}

\index{Kontsevich integral|textbf}
\index{Vassiliev invariants|textbf}

Configuration-space integrals compute knot invariants.  This is not
a metaphorical statement: the propagator
$\eta_{ij} = d\log(z_i - z_j)$ that drives the bar differential in
every preceding chapter is, upon restriction to the real
configuration space $C_n(\mathbb{R})$, the Kontsevich propagator
whose iterated integrals produce universal Vassiliev invariants.
The Kontsevich integral is therefore a sibling of the monograph's own
configuration-space calculus, built from the same differential-form
machinery on the same compactified spaces, applied to a different
geometric input (an embedded knot rather than a chiral algebra on a
curve).

The connection is structural, not accidental.  Both constructions
extract algebraic invariants by integrating logarithmic forms over
Fulton--MacPherson compactifications and reading off residues at
boundary strata.  The 4T relation governing chord diagrams is the
same Arnold relation that constrains the form space
$\Omega^*(\overline{\operatorname{Conf}}_n)$.  Weight systems ---
the linear functionals on chord diagrams that produce numerical knot
invariants --- are precisely the data carried by a metrized Lie
algebra, which is the classical shadow of an affine Kac--Moody
algebra's OPE.  This chapter makes the dictionary precise.

The propagator $\eta_{ij} = d\log(z_i - z_j)$ of the bar complex
restricts to the Kontsevich propagator on the real configuration
space, and the resulting weight systems are Vassiliev invariants
(Theorem~\ref{thm:bar-weight-systems}).

\section{Vassiliev invariants: mathematical setup}
\label{sec:vassiliev-setup}
\index{Vassiliev invariants!mathematical setup}

\begin{definition}[Vassiliev invariant]
\label{def:vassiliev-invariant}
\index{Vassiliev invariant|textbf}
A knot invariant $v\colon \{\text{oriented knots in } S^3\} \to \bC$
is a \emph{Vassiliev invariant of order~$n$} (or \emph{finite-type
invariant of type~$n$}) if it vanishes on all singular knots with
$n+1$ double points.  Equivalently, if $K_0, K_+, K_-$ denote knots
that differ only at a single crossing, the \emph{Vassiliev relation}
$v(K_\times) := v(K_+) - v(K_-)$ extends $v$ to singular knots,
and $v$ has type~$n$ if $v(K)=0$ for every singular knot $K$ with
$n+1$ transverse double points.
\end{definition}

\begin{definition}[Chord diagram]
\label{def:chord-diagram}
\index{chord diagram|textbf}
A \emph{chord diagram of order~$n$} is an oriented circle with
$n$ disjoint pairs of marked points (``chords''), considered up to
orientation-preserving diffeomorphism of the circle.  The vector
space $\cD_n$ of chord diagrams is the formal $\bC$-span of
chord diagrams of order~$n$, modulo the four-term
(\emph{4T}) relation.
\end{definition}

\begin{theorem}[Kontsevich; \ClaimStatusProvedElsewhere{} \cite{Kon94,BarNatan95}]
\label{thm:kontsevich-invariant}
\index{Kontsevich integral!universality}
The Kontsevich integral
\[
Z_K(K) \;=\;
\sum_{n=0}^\infty \;
\frac{1}{(2\pi i)^n}
\int_{\substack{t_{\min} < t_1 < \cdots < t_n < t_{\max} \\
  z_i = K(t_i),\; z'_i = K(t'_i)}}
\bigwedge_{j=1}^n
\omega^K_{j}
\;\cdot\; D_\Gamma
\]
is a universal Vassiliev invariant: every finite-type invariant
of type~$n$ factors through the degree-$n$ component of $Z_K$.
Here $\omega^K_{j} = d\log(z_j - z'_j)$ is the Kontsevich
propagator and $D_\Gamma$ is the chord diagram determined by the
pairing $(t_j, t'_j)$.
\end{theorem}

\begin{definition}[Weight system]
\label{def:weight-system}
\index{weight system|textbf}
A \emph{weight system of degree~$n$} is a linear functional
$W\colon \cD_n \to \bC$ satisfying the \emph{4T relation}.  The
\emph{Lie algebra weight system} associated to a finite-dimensional
Lie algebra~$\fg$ with non-degenerate invariant form $\kappa$ is
\[
W_\fg(\Gamma) \;=\;
\operatorname{Tr}_{\fg^{\otimes 2n}}
\!\left(
  \prod_{\substack{\text{chords}\\ (i,j)}} \kappa^{ab}_{ij}
  \;\cdot\;
  \prod_{\substack{\text{vertices}\\ v}} f^a_{bc;v}
\right),
\]
where $\kappa^{ab}$ is the inverse Killing form and $f^a_{bc}$ are
the structure constants.
\end{definition}

\section{Weight systems from the bar complex}
\label{sec:weight-systems-bar}
\index{weight system!from bar complex}

The bar differential extracts exactly the data defining a weight system.

\begin{theorem}[Bar complex weight systems; \ClaimStatusProvedHere]
\label{thm:bar-weight-systems}
For $\cA = \widehat{\fg}_k$ the affine Kac--Moody algebra at
level~$k$, the genus-$0$ bar differential defines Lie algebra
weight systems.  More precisely:
\begin{enumerate}[label=\textup{(\roman*)}]
\item For each chord diagram $\Gamma$ of degree~$n$, the bar
  complex element
  $\alpha_\Gamma \in \barB^n(\widehat{\fg}_k)$
  defined by
  \[
  \alpha_\Gamma \;=\;
  \bigwedge_{\substack{\text{chords}\\(i,j) \text{ of } \Gamma}}
  \eta_{ij}
  \;\otimes\;
  \bigotimes_{v=1}^{2n} J^{a_v}
  \]
  where $\eta_{ij} = d\log(z_i - z_j)$ is the FM propagator on
  $\overline{C}_{2n}(\bP^1)$ and $J^{a_v}$ is the Kac--Moody
  current, satisfies
  \[
  \langle d^{(0)} \alpha_\Gamma, \mathbf{1} \rangle
  \;=\;
  k^n \cdot W_\fg(\Gamma).
  \]
\item The 4T relation for weight systems is a consequence of the
  Arnold relations for the propagator:
  \[
  \eta_{ij} \wedge \eta_{jk}
  + \eta_{jk} \wedge \eta_{ki}
  + \eta_{ki} \wedge \eta_{ij} = 0.
  \]
\end{enumerate}
\end{theorem}

\begin{proof}
Part~(i) follows from the bar differential formula
(Definition~\ref{def:bar-differential-complete}): at genus~$0$, the
bar differential is $d = \sum_{i<j} \Res_{D_{ij}}(\eta_{ij}
\wedge -)$, where the residue extracts the OPE coefficient.  For
the Kac--Moody current $J^a(z_i) J^b(z_j)
\sim k\kappa^{ab}/(z_i-z_j)^2 + f^{ab}_c J^c/(z_i-z_j)$, the
double-pole residue gives the Killing contraction $k\kappa^{ab}$
and the simple-pole residue gives the structure constants
$f^{ab}_c$.  The pairing with the vacuum $\mathbf{1}$ projects onto
the scalar (double-pole) contribution, yielding exactly
$k^n W_\fg(\Gamma)$.

Part~(ii): the 4T relation for weight systems states that for four
points $i,j,k,l$ on the circle, the sum of the three resolutions of
the crossing pattern vanishes.  This follows from the Arnold
relation on $\overline{C}_4(\bP^1)$, which ensures that the bar
differential $d^2 = 0$ at degree~$2$: the Arnold relation
$\eta_{ij}\eta_{jk} + \eta_{jk}\eta_{ki}
+ \eta_{ki}\eta_{ij} = 0$ implies the same linear combination
of weight system evaluations vanishes.  (The Arnold relation is
the OS relation on three generators; the 4T relation is its
image under the OPE pairing.)
\end{proof}

This theorem refines Proposition~\ref{prop:vassiliev-genus0} by
making the chord diagram correspondence explicit.

\section{The Kontsevich propagator from holomorphic restriction}
\label{sec:kontsevich-propagator}
\index{Kontsevich integral!propagator restriction}

The holomorphic propagator restricts to the Kontsevich propagator on the real locus.

\begin{proposition}[Propagator restriction; \ClaimStatusProvedHere]
\label{prop:propagator-restriction}
Let $S^1 = \{|z| = 1\} \subset \bC$ be the unit circle, parametrized
by $z = e^{i\theta}$.  The FM propagator $\eta_{ij} = d\log(z_i-z_j)$
restricts on $C_n(S^1) \hookrightarrow C_n(\bC)$ as
\[
\eta_{ij}\big|_{S^1}
\;=\;
\frac{d(z_i - z_j)}{z_i - z_j}\bigg|_{z_k = e^{i\theta_k}}
\;=\;
\frac{i\, e^{i\theta_i}\, d\theta_i - i\, e^{i\theta_j}\, d\theta_j}
     {e^{i\theta_i} - e^{i\theta_j}}.
\]
Taking the real part and using
$e^{i\theta_i}/(e^{i\theta_i}-e^{i\theta_j})
= \tfrac{1}{2} + \tfrac{i}{2}\cot((\theta_i-\theta_j)/2)$:
\[
\Re(\eta_{ij})\big|_{S^1}
\;=\;
\frac{1}{2}\, d\theta_i
+ \frac{1}{2}\cot\!\left(\frac{\theta_i-\theta_j}{2}\right)
  \frac{d\theta_i + d\theta_j}{2}
\;=\;
\frac{1}{2\pi}\, d\arg(e^{i\theta_i} - e^{i\theta_j})
\;+\; \text{\emph{exact}},
\]
where the exact term does not contribute to closed integrals
over configuration spaces.  Hence:
\[
\int_{C_n(S^1)} \prod_{(i,j) \in \Gamma}
\Re(\eta_{ij})
\;=\;
\int_{C_n(S^1)} \prod_{(i,j) \in \Gamma}
\eta^K_{ij}
\]
where $\eta^K_{ij} = \frac{1}{2\pi}\, d\arg(t_i - t_j)$ is
the Kontsevich propagator.
\end{proposition}

\begin{proof}
Write $z_k = e^{i\theta_k}$, so $dz_k = ie^{i\theta_k}\, d\theta_k$.
Then
\[
\frac{dz_i}{z_i - z_j}
= \frac{i\, e^{i\theta_i}\, d\theta_i}
       {e^{i\theta_i} - e^{i\theta_j}}.
\]
Factor out $e^{i(\theta_i+\theta_j)/2}$:
\[
\frac{e^{i\theta_i}}{e^{i\theta_i}-e^{i\theta_j}}
= \frac{e^{i(\theta_i-\theta_j)/2}}
       {e^{i(\theta_i-\theta_j)/2} - e^{-i(\theta_i-\theta_j)/2}}
= \frac{1}{1 - e^{-i(\theta_i-\theta_j)}}
= \frac{1}{2} + \frac{i}{2}\cot\!\frac{\theta_i-\theta_j}{2}.
\]
The real part of $\eta_{ij}|_{S^1} = d\log(z_i-z_j)|_{S^1}$
splits as
$\Re(\eta_{ij}) = d\log|z_i-z_j| = d\log|e^{i\theta_i}-e^{i\theta_j}|
= d\log(2|\sin((\theta_i-\theta_j)/2)|)$.  For closed chord diagram
integrals, the exact form $d\log|\cdot|$ integrates to zero on the
compact fibers.  What remains is the angular part,
$\Im(\eta_{ij})|_{S^1} = d\arg(z_i-z_j)$, which is the Kontsevich
propagator (up to normalization $1/(2\pi)$).
\end{proof}

\begin{remark}[Imaginary part]\label{rem:imaginary-part}
In the Kontsevich integral, the propagator
$\eta^K_{ij} = (1/(2\pi i))\, d\log(z_i - z_j)$
uses the full holomorphic logarithm, not just the real part.
The distinction is irrelevant for the leading-order weight systems
but matters for the framing correction.  More precisely:
the real part $d\log|z_i-z_j|$ contributes the ``anomalous'' part
of the Kontsevich integral (the writhe/framing term), while the
imaginary part $d\arg(z_i-z_j)$ contributes the ``invariant'' part.
The bar complex naturally produces both contributions via the
full holomorphic propagator.
\end{remark}


\section{The Drinfeld associator from bar-cobar monodromy}
\label{sec:drinfeld-associator}
\index{Drinfeld associator|textbf}
\index{KZ equation!from bar complex}

The bar complex on~$C_3(\bC)$ carries a flat connection whose monodromy is the Drinfeld associator.

\begin{proposition}[KZ connection from bar complex; \ClaimStatusProvedHere]
\label{prop:kz-from-bar}
\index{KZ equation!from bar differential}
For $\cA = \widehat{\fg}_k$, the bar differential restricted to
$C_n(\bC) \subset \overline{C}_n(\bP^1)$ defines a flat connection
on the trivial bundle with fiber $U(\fg)^{\otimes n}$:
\begin{equation}\label{eq:kz-connection}
\nabla^{\mathrm{KZ}}
\;=\;
d \;-\; \frac{1}{k+h^\vee}
\sum_{1 \leq i < j \leq n}
\frac{\Omega_{ij}}{z_i - z_j}\, d(z_i - z_j),
\end{equation}
where $\Omega_{ij} = \sum_a J^a_i J^a_j$ is the Casimir element
acting on the $i$-th and $j$-th tensor factors.  This is the
Knizhnik--Zamolodchikov connection.
\end{proposition}

\begin{proof}
The bar differential for $\widehat{\fg}_k$ at genus~$0$
(Definition~\ref{def:bar-differential-complete}) is
$d = \sum_{i<j} \Res_{D_{ij}} (\eta_{ij} \wedge -)$,
where $\eta_{ij} = d\log(z_i-z_j)$ and $\Res_{D_{ij}}$
extracts the simple-pole OPE coefficient.  The simple-pole term in
the OPE $J^a(z_i) J^b(z_j) \sim \cdots + f^{ab}_c J^c/(z_i-z_j)$
is the adjoint action.  Contracting with the Killing form gives
$\Omega_{ij}/(z_i-z_j)$.  The factor $1/(k+h^\vee)$ arises from
the Sugawara normalization: the OPE is formulated at level~$k$,
and the KZ connection uses the normalized Casimir
$\Omega/(k+h^\vee)$.  (At the critical level $k = -h^\vee$,
the Sugawara construction is undefined and the KZ connection
degenerates; the correct replacement is the Feigin--Frenkel
connection on opers.)

Flatness $(\nabla^{\mathrm{KZ}})^2 = 0$ follows from $d^2 = 0$
on the genus-$0$ bar complex (the bar differential restricted to
the open stratum $C_n(\bC)$ has no curvature since $\kappa$ arises
only at genus~$\geq 1$).  Explicitly, the flatness condition is
the infinitesimal braid relation
$[\Omega_{ij}, \Omega_{ik} + \Omega_{jk}] = 0$ for $i,j,k$ distinct,
which follows from the Jacobi identity for~$\fg$.
\end{proof}

\begin{theorem}[Drinfeld associator from bar-cobar; \ClaimStatusProvedHere]
\label{thm:drinfeld-associator-bar}
\index{Drinfeld associator!from bar monodromy}
The monodromy of the KZ connection~\eqref{eq:kz-connection}
on $C_3(\bC)$ around the boundary divisors of
$\overline{C}_3(\bP^1)$ produces the Drinfeld associator
$\Phi_{\mathrm{KZ}} \in U(\fg)^{\otimes 3}[[\hbar]]$
(where $\hbar = 2\pi i/(k+h^\vee)$).  In bar-cobar terms:
\begin{enumerate}[label=\textup{(\roman*)}]
\item The parallel transport of the bar complex
  $\barB(\widehat{\fg}_k)$ around the loop
  $\{z_1 = 0,\, z_2 = 1,\, z_3 = t\}$ for $t \in (0,1)$
  with $t \to 0$ and $t \to 1$ as the two boundary degenerations
  gives the regularized holonomy
  $\Phi_{\mathrm{KZ}}(t_{12}, t_{23})$
  where $t_{ij} = \Omega_{ij}/(k+h^\vee)$.
\item The pentagon relation for the associator
  is a consequence of the coherence of the bar complex
  on $C_4(\bC)$: the five boundary strata of
  $\overline{\mathcal{M}}_{0,5} \cong \operatorname{Bl}_4(\bP^2)$
  give five ways to degenerate four-point configurations,
  and the associator intertwines them.
\item The hexagon relation comes from the braid monodromy
  on $C_3(\bC)$: the half-twist $\sigma_1$ exchanges
  $z_1$ and $z_2$, and
  $\Phi_{\mathrm{KZ}} \circ \sigma_1 \circ \Phi_{\mathrm{KZ}}^{-1}
  = R_{12}$
  where $R_{12} = e^{\pi i\, t_{12}}$ is the $R$-matrix.
\end{enumerate}
\end{theorem}

\begin{proof}
The Drinfeld--Kohno theorem (\ClaimStatusProvedElsewhere,
\cite{Drinfeld90,Kohno87}) establishes that the KZ connection
produces the associator; our contribution is the identification
of the KZ connection with the bar complex structure.

Part (i): the bar differential on $C_3(\bC)$ is
$d = t_{12}\, \eta_{12} + t_{13}\, \eta_{13} + t_{23}\, \eta_{23}$
where $\eta_{ij} = d\log(z_i - z_j)$ and
$t_{ij} = \Omega_{ij}/(k+h^\vee)$.  The parallel transport from
$t = \varepsilon$ to $t = 1 - \varepsilon$ is
$\Phi_{\mathrm{KZ}}(t_{12}, t_{23})
= 1 + \sum_{n \geq 1} \int_{0<s_1<\cdots<s_n<1}
\omega(s_1) \cdots \omega(s_n)$
where $\omega(s) = t_{12}\, ds/(s) + t_{23}\, ds/(s-1)$, regularized
at $s = 0$ and $s = 1$.  This is the iterated integral formula for
the Drinfeld associator.

Part (ii): on $C_4(\bC)$, the bar complex structure requires
coherence of all four-point degenerations.  The boundary of
$\overline{\mathcal{M}}_{0,5}$ has five components (the five
Ptolemy relations), and the parallel transports around these
components compose to give the pentagon identity
$\Phi_{234} \cdot \Phi_{124} = \Phi_{123} \cdot \Phi_{134} \cdot \Phi_{234}$
(Drinfeld's equation).

Part (iii): the half-twist $\sigma_1$ on $C_3(\bC)$ gives
the braid generator, and the monodromy of the KZ connection
under this half-twist is $R_{12} = e^{\pi i\, t_{12}}$
(the universal $R$-matrix at level~$k$).  The hexagon relation
expresses compatibility of the associator with the braiding.
\end{proof}

\begin{remark}[Bar-cobar as quasi-triangular structure]
\label{rem:bar-quasi-triangular}
\index{quasi-triangular!from bar complex}
Theorem~\ref{thm:drinfeld-associator-bar} shows that the bar
complex of $\widehat{\fg}_k$ naturally produces the data of a
quasi-triangular quasi-Hopf algebra: the associator~$\Phi$ from
3-point monodromy, the $R$-matrix from 2-point braiding, and the
pentagon/hexagon from 4/3-point coherence.  This is
\emph{not} an analogy: the bar differential \emph{is} the
KZ connection, and the bar-cobar adjunction \emph{is} the
Reshetikhin--Turaev construction at the chain level.
The chiral operad on $\{C_n(\bC)\}_{n \geq 2}$ encodes the
full quasi-triangular structure; Koszul duality (bar-cobar
inversion) exchanges the braiding $R$ with its
inverse~$R^{-1}$ --- precisely the passage from Yangian to
quantum group that the $\Eone$-chiral Koszul duality theorem
(Theorem~\ref{thm:e1-chiral-koszul-duality}) makes rigorous.
\end{remark}


\section{Higher genus: from bar complex to loop expansion}
\label{sec:bar-to-loop}
\index{Kontsevich integral!loop expansion}
\noindent\emph{Status split.} The graph-complex identification is
\ClaimStatusProvedElsewhere/\ClaimStatusProvedHere as tagged below; the
holomorphic-to-topological comparison for $g \geq 1$ remains
\ClaimStatusConjectured and is isolated as the only missing bridge.

The genus-$0$ results extend conjecturally to the full genus tower:
the algebraic bar-complex structure extends to all genera (proved);
the holomorphic-to-topological propagator comparison is the conjectural step.

\begin{proposition}[Feynman transform produces graph complex;
\ClaimStatusProvedElsewhere]
\label{prop:feynman-graph-complex}
\index{Feynman transform!graph complex}
The Feynman transform of the modular operad $\mathrm{Com}$
recovers the Kontsevich graph complex computing Vassiliev
invariants.  Specifically:
\begin{enumerate}[label=\textup{(\roman*)}]
\item (\ClaimStatusProvedElsewhere) By Getzler--Kapranov~\cite{GeK98},
  the Feynman transform
  $\mathrm{FT}(\mathrm{Com})$ is the graph complex $\mathrm{GC}$
  whose cohomology computes the rational homotopy type of the
  space of long knots.
\item (\ClaimStatusProvedHere) Our identification
  $\barB^{\mathrm{full}}(\cA)
  \simeq \mathrm{FT}_{\mathcal{M}\mathrm{od}}(\cA)$
  (Theorem~\ref{thm:prism-higher-genus}) extends this to Koszul
  chiral algebras: the full bar complex is a refinement of the
  graph complex with coefficients in the OPE data of~$\cA$.
  This identification is \emph{algebraic}: it matches the combinatorial
  structures (graph types, differentials, genus grading) on both sides.
\item (\ClaimStatusConjectured) The passage from holomorphic propagators on complex curves
  to the topological propagators of the Kontsevich integral is
  established at genus~$0$ (Proposition~\ref{prop:propagator-restriction})
  but remains conjectural at genus $g \geq 1$; see
  Remark~\ref{rem:holomorphic-to-topological} and
  Conjecture~\ref{conj:vassiliev-bar}.
\end{enumerate}
\end{proposition}

\begin{proof}
Part (i) is Getzler--Kapranov \cite{GeK98}, Theorem~5.4.
Part~(ii): Theorem~\ref{thm:prism-higher-genus} identifies the
genus-$g$ bar complex with the genus-$g$ component of the
Feynman transform $\mathrm{FT}_{\mathcal{M}\mathrm{od}}(\cA)$.
At $\cA = \mathrm{Com}$ (the commutative operad, viewed as a
trivial chiral algebra), the OPE data is trivial and
$\mathrm{FT}_{\mathcal{M}\mathrm{od}}(\mathrm{Com}) =
\mathrm{FT}(\mathrm{Com}) = \mathrm{GC}$.
\end{proof}

\begin{remark}[The gap: holomorphic to topological]
\label{rem:holomorphic-to-topological}
\index{Kontsevich integral!holomorphic vs topological}
The remaining gap is the passage from holomorphic to topological:
\begin{enumerate}
\item Our propagators are holomorphic forms on complex curves;
  the Kontsevich propagator is a real form on~$S^1$.
\item Proposition~\ref{prop:propagator-restriction} handles the
  genus-$0$ case: the holomorphic propagator restricts to the
  Kontsevich propagator on $S^1 \subset \bP^1$.
\item At genus $g \geq 1$, the holomorphic propagator
  $\omega(z,w) = d_z \log E(z,w)$ (where $E$ is the prime form)
  should restrict to the appropriate real propagator on a real
  cycle $S^1 \hookrightarrow \Sigma_g$, but this restriction
  requires Chern--Simons gauge equivalence
  (Cattaneo--Mnev~\cite{CattaneoMnev10}), not merely
  pointwise restriction.
\end{enumerate}
This gap is the subject of Conjecture~\ref{conj:vassiliev-bar}
in \S\ref{subsec:vassiliev}.
\end{remark}


\section{Chern--Simons theory as the bridge}
\label{sec:cs-bridge}
\index{Chern--Simons!as bridge to Vassiliev}
Chern--Simons gauge theory is the proposed bridge from holomorphic bar
complexes to topological knot invariants.

\begin{conjecture}[CS factorization homology; \ClaimStatusConjectured]
\label{conj:cs-factorization}
\index{Chern--Simons!factorization homology}
For a compact oriented $3$-manifold~$M$ with
$\partial M = \Sigma$ a Riemann surface, the perturbative
Chern--Simons partition function at level~$k$ satisfies:
\begin{enumerate}[label=\textup{(\roman*)}]
\item \textup{(Boundary-to-bulk)}\quad
  The algebra of observables $\mathrm{Obs}_{\mathrm{CS}}(M)$
  is the factorization homology
  $\int_M \barB(\widehat{\fg}_k)_\Sigma$,
  where $\barB(\widehat{\fg}_k)_\Sigma$ denotes the bar complex
  of $\widehat{\fg}_k$ viewed as a factorization algebra on~$M$
  via the holomorphic-topological twist along the
  surface~$\Sigma = \partial M$.
\item \textup{(Knot invariants)}\quad
  For $M = S^3$ and a knot $K \subset S^3$ colored by a
  representation $V$ of~$\fg$, the expectation value
  $\langle W_K(V) \rangle_{\mathrm{CS}}$ is computed by the
  bar complex with module insertion:
  $\langle W_K(V) \rangle = \chi(\barB(\widehat{\fg}_k, V_K))$
  where $V_K$ is the $\widehat{\fg}_k$-module supported on~$K$.
\end{enumerate}
\end{conjecture}

\begin{remark}[Scope]\label{rem:cs-scope}
This conjecture connects three bodies of work:
(a)~our bar complex on curves (proved);
(b)~Costello--Gwilliam's perturbative quantization of CS
(\cite{CG17}, constructing the factorization algebra);
(c)~Axelrod--Singer's perturbative CS (\cite{Kon94}),
computing knot invariants via configuration space integrals.
The identification (i) requires showing that Costello--Gwilliam's
factorization algebra for CS, restricted to the boundary, agrees
with our bar complex.  This is the content of Programme~VI-a
(\S\ref{subsec:anomaly-koszul}).

\emph{Homotopy template} (Convention~\ref{conv:hms-levels}): Type~VII (physics-dictionary).
The conjecture treats the bar complex of $\widehat{\fg}_k$ as the
boundary-side algebraic shadow of the perturbative Chern--Simons
factorization algebra, with knot-invariant comparison mediated by bar
cohomology with module insertions.
(Contributing to Conjecture~\ref{conj:master-bv-brst}.)
\end{remark}


\section{Summary}
\label{sec:kontsevich-summary}

\noindent
\textbf{What is proved.}
\begin{enumerate}
\item Bar complex weight systems equal Lie algebra weight systems
  (Theorem~\ref{thm:bar-weight-systems}).
\item The holomorphic propagator restricts to the Kontsevich
  propagator on $S^1 \subset \bP^1$
  (Proposition~\ref{prop:propagator-restriction}).
\item The bar differential on $C_n(\bC)$ is the KZ connection
  (Proposition~\ref{prop:kz-from-bar}).
\item The Drinfeld associator arises from bar-cobar monodromy on
  $C_3(\bC)$ (Theorem~\ref{thm:drinfeld-associator-bar}).
\item The Feynman transform of $\mathrm{Com}$ is the graph complex
  (Proposition~\ref{prop:feynman-graph-complex}).
\end{enumerate}

\noindent
\textbf{What remains.}
The full identification of the genus tower with the loop expansion
of the Kontsevich integral requires the Chern--Simons bridge at
higher genus (Conjecture~\ref{conj:vassiliev-bar},
Conjecture~\ref{conj:cs-factorization}).  The genus-$0$ case is
established; the gap is the passage from holomorphic to topological
Feynman transforms at $g \geq 1$.
